\\documentclass[11pt]{article}
\\usepackage[T1]{fontenc}
\\usepackage[utf8]{inputenc}
\\usepackage{geometry}
\\usepackage{xcolor}
\\usepackage{fvextra}
\\usepackage{hyperref}
\\geometry{margin=1in}
\\DefineVerbatimEnvironment{CodeBlock}{Verbatim}{breaklines,breakanywhere,fontsize=\\small}

\begin{document}
\section*{core/metric\\_gradients.py}
\textbf{Source Path}: \texttt{core/metric\\_gradients.py}\par
\bigskip
\begin{CodeBlock}
"""
metric_gradients.py - Semantic Tension Mapping via Metric Learning Gradients

SEPARATE LANE from embedding-based approach.

THE INSIGHT:
Instead of classifying (NLI logits) or clustering (embeddings), we measure
the GRADIENT of geometric alignment. This tells us the "force vector" -
how an article would need to change to align with a pure concept.

THE PHYSICS:
- Old way: Stand on probability hill, ask "which way is up?"
  Problem: Hill is flat at 99% confidence. Gradients vanish.
  
- New way: Place magnetic anchors (pure concepts), measure pull on article.
  Fix: Even far articles have pull vectors. No confidence plateau.

CRITICAL: Uses Bi-Encoder (NOT Cross-Encoder) to prevent attention leakage.
- Pass 1: Encode Article alone -> Article_Vector
- Pass 2: Encode Anchor alone -> Anchor_Vector  
- Compute: cosine(Article_Vector, Anchor_Vector)
- Backprop: gradient tells us "rotation toward anchor"

WHY GRADIENTS LIVE IN TANGENT SPACE:
Grad(cosine) is orthogonal to the vector itself.
This is a FEATURE: filters out magnitude noise, isolates direction (framing).

THESIS REFRAME (if this works):
"Semantic Tension Mapping: We treat embedding space as a force field.
Different framing concepts exert different pulls. The angle between pull
vectors reveals whether an article is semantically ambiguous (parallel pulls)
or locked into a frame (opposing pulls)."

Author: Belief Transformer Project
"""

from __future__ import annotations

import hashlib
import json
import re
import datetime
from pathlib import Path
from typing import Dict, List, Optional, Tuple, Any, Union
from dataclasses import dataclass, asdict, field

import torch
import torch.nn as nn
import torch.nn.functional as F
import numpy as np


# =============================================================================
# PARAGRAPH UTILITIES (mirrored from nli_extraction.py for consistency)
# =============================================================================

def _split_into_paragraphs(text: str, max_paragraphs: int) -> List[str]:
    """Split text on double-newline boundaries."""
    if not isinstance(text, str):
        return []
    t = text.strip()
    if not t:
        return []
    chunks = re.split(r"\n\s*\n+", t)
    paras = [c.strip() for c in chunks if c and c.strip()]
    return paras[:max_paragraphs] if max_paragraphs and max_paragraphs > 0 else paras


def _normalize_weights(weights: List[float], n: int) -> List[float]:
    """Normalize weight vector to sum to 1.0, handling variable paragraph counts."""
    if n <= 0:
        return []
    w = weights[:n] if weights else []
    if len(w) < n:
        w = w + [w[-1] if w else 1.0] * (n - len(w))
    s = float(sum(w))
    if s <= 0:
        return [1.0 / n] * n
    return [float(x) / s for x in w]


# =============================================================================
# TRACK 3: GEOMETRIC STABILITY (Modern System Required)
# =============================================================================
# Added to support complete_pipeline.py calls without breaking legacy code.

def compute_metric_gradients(
    sweep_results: Dict[str, Any],
    alphas: List[float],
) -> Dict[str, Any]:
    """
    Compute geometric stability gradients from an alpha sweep.
    
    Equation: Tension(α) ≈ || Gram(α+Δ) - Gram(α) ||_F / Δα
    """
    wavelengths = sweep_results.get("wavelength", {})
    
    if not wavelengths:
        # Fallback if no sweep data
        return _compute_gradients_from_gram(sweep_results, alphas)

    gradients = []
    intervals = []
    sorted_alphas = sorted(alphas)
    
    for i in range(len(sorted_alphas) - 1):
        a1 = sorted_alphas[i]
        a2 = sorted_alphas[i+1]
        
        # Robust key retrieval (handles different formatting)
        key_candidates = [
            f"{a1}_to_{a2}", 
            f"{a2}_to_{a1}",
            f"alpha_{a1}_to_{a2}"
        ]
        
        data = None
        for k in key_candidates:
            if k in wavelengths:
                data = wavelengths[k]
                break
            
        if data is None:
            continue
            
        # Extract energy
        if isinstance(data, dict):
            energy = data.get("energy", data.get("frobenius_norm", 0.0))
        else:
            energy = float(data)
            
        delta_alpha = abs(a2 - a1)
        gradient = energy / (delta_alpha + 1e-9)
        
        gradients.append(float(gradient))
        intervals.append(f"{a1}-{a2}")

    if not gradients:
        return {"mean_tension": 0.0, "stability_score": 1.0, "gradients": []}

    # Aggregate Stats
    gradients_arr = np.array(gradients)
    mean_tension = float(np.mean(gradients_arr))
    max_tension = float(np.max(gradients_arr))
    std_tension = float(np.std(gradients_arr))
    
    # Critical Phase Transition
    max_idx = int(np.argmax(gradients_arr))
    critical_interval = intervals[max_idx]
    
    # Phase transitions
    threshold = mean_tension + std_tension
    phase_transitions = [intervals[i] for i, g in enumerate(gradients) if g > threshold]

    # Stability Score: 1.0 = Rigid, 0.0 = Fluid
    stability_score = 1.0 / (1.0 + mean_tension)

    return {
        "gradients": gradients,
        "intervals": intervals,
        "mean_tension": mean_tension,
        "max_tension": max_tension,
        "std_tension": std_tension,
        "critical_interval": critical_interval,
        "phase_transitions": phase_transitions,
        "stability_score": stability_score,
        "n_steps": len(gradients)
    }

def _compute_gradients_from_gram(sweep_results, alphas):
    """Fallback for when pre-computed wavelengths are missing."""
    gradients = []
    intervals = []
    sorted_alphas = sorted(alphas)

    for i in range(len(sorted_alphas) - 1):
        a1 = sorted_alphas[i]
        a2 = sorted_alphas[i + 1]
        
        g1 = sweep_results.get(f"alpha_{a1}", {}).get("gram_matrix")
        g2 = sweep_results.get(f"alpha_{a2}", {}).get("gram_matrix")

        if g1 is None or g2 is None: continue

        if hasattr(g1, 'cpu'): g1 = g1.cpu().numpy()
        if hasattr(g2, 'cpu'): g2 = g2.cpu().numpy()

        diff = g2 - g1
        energy = float(np.linalg.norm(diff, ord='fro'))
        gradient = energy / (abs(a2 - a1) + 1e-9)

        gradients.append(gradient)
        intervals.append(f"{a1}-{a2}")

    if not gradients:
        return {"mean_tension": 0.0, "stability_score": 1.0}

    gradients_arr = np.array(gradients)
    mean_tension = float(np.mean(gradients_arr))
    stability_score = 1.0 / (1.0 + mean_tension)

    return {
        "gradients": gradients,
        "intervals": intervals,
        "mean_tension": mean_tension,
        "stability_score": stability_score,
        "method": "gram_fallback"
    }

# Alias for backward compatibility
compute_alpha_gradients = compute_metric_gradients


@dataclass
class AlphaStabilityConfig:
    """Configuration for alpha-sweep stability analysis (Track 3)."""
    alphas: List[float] = field(default_factory=lambda: [0.1, 0.5, 1.0, 5.0, 20.0])
    n_dirichlet_samples: int = 50
    rks_output_dim: int = 2048
    crn_enabled: bool = True
    crn_seed: int = 12345
    device: str = "cuda"
    def to_dict(self) -> Dict: return asdict(self)


class AlphaStabilityAnalyzer:
    """Analyze geometric stability across Dirichlet alpha values (Track 3)."""
    
    def __init__(self, config: AlphaStabilityConfig = None):
        self.config = config or AlphaStabilityConfig()
        self.device = torch.device(self.config.device if torch.cuda.is_available() else 'cpu')
        print(f"[AlphaStability] Initialized with {len(self.config.alphas)} alpha values")
        
    def compute_gram_matrix(self, embeddings: torch.Tensor, normalize: bool = True) -> torch.Tensor:
        if normalize: embeddings = F.normalize(embeddings, p=2, dim=-1)
        return torch.mm(embeddings, embeddings.t())

    def run_alpha_sweep(self, cls_per_bot: torch.Tensor) -> Dict[str, Any]:
        N, n_bots, hidden = cls_per_bot.shape
        cls_per_bot = cls_per_bot.to(self.device)
        
        if self.config.crn_enabled:
            torch.manual_seed(self.config.crn_seed)
            
        results = {}
        gram_matrices = {}
        
        for alpha in self.config.alphas:
            alpha_vec = torch.full((n_bots,), alpha, device=self.device)
            weights = torch.distributions.Dirichlet(alpha_vec).sample((self.config.n_dirichlet_samples,))
            
            # Fuse
            weights_exp = weights.unsqueeze(0).unsqueeze(-1)
            cls_exp = cls_per_bot.unsqueeze(1)
            fused = (weights_exp * cls_exp).sum(dim=2)
            fused_mean = fused.mean(dim=1)
            fused_std = fused.std(dim=1).mean().item()
            
            gram = self.compute_gram_matrix(fused_mean)
            gram_matrices[alpha] = gram
            
            results[f"alpha_{alpha}"] = {
                "fused": fused_mean.cpu(),
                "fused_std": fused_std,
                "gram_matrix": gram.cpu(),
            }
            
        wavelengths = {}
        sorted_alphas = sorted(self.config.alphas)
        for i in range(len(sorted_alphas) - 1):
            a1, a2 = sorted_alphas[i], sorted_alphas[i+1]
            diff = gram_matrices[a2] - gram_matrices[a1]
            energy = float(torch.norm(diff, p='fro').item())
            wavelengths[f"{a1}_to_{a2}"] = {"energy": energy, "delta_alpha": a2 - a1}
            
        results["wavelength"] = wavelengths
        results["alphas"] = self.config.alphas
        return results

    def analyze_corpus(self, articles: List[Dict], cls_per_bot: torch.Tensor) -> Dict[str, Any]:
        print(f"\n[AlphaStability] Analyzing {len(articles)} articles...")
        sweep_results = self.run_alpha_sweep(cls_per_bot)
        gradients = compute_metric_gradients(sweep_results, self.config.alphas)
        
        return {
            "sweep_results": sweep_results,
            "metric_gradients": gradients,
            "config": self.config.to_dict(),
            "n_articles": len(articles)
        }


# =============================================================================
# CONFIGURATION
# =============================================================================

@dataclass
class MetricGradientConfig:
    """Configuration for metric gradient extraction."""
    
    # Model settings
    model_name: str = "microsoft/deberta-v3-base"
    device: str = "cuda"
    
    # Anchor definitions (pure concepts)
    anchors: Dict[str, str] = field(default_factory=lambda: {
        'victim': "This text describes victims and suffering.",
        'aggressor': "This text describes aggression and violence.",
        'neutral': "This text is a neutral factual report.",
        'emotional': "This text is emotionally charged.",
        'humanitarian': "This text focuses on humanitarian concerns.",
        'security': "This text focuses on security threats.",
        'conflict': "This text describes armed conflict.",
        'peace': "This text describes peace and resolution.",
    })
    
    # Extraction settings
    normalize_gradients: bool = True
    gradient_layer: str = 'last'  # 'last', 'second_last', 'mean'
    pooling: str = 'mean'  # 'cls', 'mean', 'last_token'
    
    # Paragraph awareness (matches nli_extraction.py hard invariant)
    paragraph_aware: bool = True
    paragraph_weights: List[float] = field(default_factory=lambda: [0.5, 0.3, 0.2])

    # Analysis settings
    compute_tension_matrix: bool = True  # Compute all pairwise anchor tensions
    
    def to_dict(self) -> Dict:
        return asdict(self)


# =============================================================================
# METRIC GRADIENT EXTRACTOR (Bi-Encoder)
# =============================================================================

class MetricGradientExtractor:
    """
    Extract metric learning gradients for semantic tension mapping.
    
    Uses Bi-Encoder approach to prevent attention leakage:
    - Article encoded separately (no hypothesis contamination)
    - Anchor encoded separately (pure concept)
    - Gradient of cosine similarity computed
    
    The gradient = "force vector" toward the anchor concept.
    """
    
    def __init__(self, config: MetricGradientConfig = None):
        self.config = config or MetricGradientConfig()
        self.device = torch.device(self.config.device if torch.cuda.is_available() else 'cpu')
        
        # Lazy load model
        self._model = None
        self._tokenizer = None
        
        # Cache anchor vectors (computed once)
        self._anchor_cache: Dict[str, torch.Tensor] = {}
        
        print(f"[MetricGradient] Initialized with {len(self.config.anchors)} anchors")
        print(f"[MetricGradient] Device: {self.device}")
    
    @property
    def model(self):
        if self._model is None:
            self._load_model()
        return self._model
    
    @property
    def tokenizer(self):
        if self._tokenizer is None:
            self._load_model()
        return self._tokenizer
    
    def _load_model(self):
        """Lazy load the transformer model."""
        from transformers import AutoModel, AutoTokenizer
        
        print(f"[MetricGradient] Loading {self.config.model_name}...")
        self._tokenizer = AutoTokenizer.from_pretrained(self.config.model_name)
        self._model = AutoModel.from_pretrained(self.config.model_name)
        self._model.to(self.device)
        self._model.eval()
        print(f"[MetricGradient] Model loaded")
    
    def _get_anchor_vector(self, anchor_name: str) -> torch.Tensor:
        """Get cached anchor vector (pure concept embedding)."""
        if anchor_name not in self._anchor_cache:
            anchor_text = self.config.anchors[anchor_name]
            
            with torch.no_grad():
                inputs = self.tokenizer(
                    anchor_text, 
                    return_tensors='pt', 
                    padding=True, 
                    truncation=True,
                    max_length=128,
                )
                inputs = {k: v.to(self.device) for k, v in inputs.items()}
                
                output = self.model(**inputs)
                
                # Pool
                if self.config.pooling == 'cls':
                    anchor_vec = output.last_hidden_state[:, 0, :]
                elif self.config.pooling == 'mean':
                    mask = inputs['attention_mask'].unsqueeze(-1).float()
                    anchor_vec = (output.last_hidden_state * mask).sum(dim=1) / mask.sum(dim=1)
                else:
                    anchor_vec = output.last_hidden_state[:, -1, :]
                
                # Normalize to unit sphere
                anchor_vec = F.normalize(anchor_vec, p=2, dim=1)
                
                self._anchor_cache[anchor_name] = anchor_vec
        
        return self._anchor_cache[anchor_name]
    
    def get_metric_gradient(
        self,
        article_text: str,
        anchor_name: str,
    ) -> torch.Tensor:
        """
        Compute gradient of article embedding moving toward anchor.

        Uses Bi-Encoder: Article encoded alone, no attention leakage.
        When paragraph_aware=True, splits text into paragraphs and
        accumulates weighted gradients (matching nli_extraction.py invariant).

        Args:
            article_text: The article to analyze
            anchor_name: Name of anchor concept (key in config.anchors)

        Returns:
            gradient: [hidden_dim] force vector toward anchor
        """
        if self.config.paragraph_aware:
            return self._get_gradient_paragraph_aware(article_text, anchor_name)
        return self._get_gradient_single(article_text, anchor_name)

    def _get_gradient_paragraph_aware(
        self,
        article_text: str,
        anchor_name: str,
    ) -> torch.Tensor:
        """
        Paragraph-weighted gradient accumulation.

        ∇_final = Σ_i w_i · ∇_paragraph_i

        This ensures the Force field (Track 1.5) matches the structural
        weighting of the Embedding field (Track 2), so Phase Space fusion
        (Track 4) operates on geometrically consistent inputs.
        """
        weights = self.config.paragraph_weights
        paras = _split_into_paragraphs(article_text, len(weights))

        if not paras:
            paras = [article_text.strip()]

        w = _normalize_weights(weights, len(paras))

        gradient_acc = None
        for para, weight in zip(paras, w):
            para_grad = self._get_gradient_single(para, anchor_name)
            if gradient_acc is None:
                gradient_acc = para_grad * weight
            else:
                gradient_acc = gradient_acc + para_grad * weight

        # Re-normalize the accumulated gradient (weighted sum may shrink magnitude)
        if self.config.normalize_gradients:
            gradient_acc = F.normalize(gradient_acc.unsqueeze(0), p=2, dim=1).squeeze(0)

        return gradient_acc

    def _get_gradient_single(
        self,
        text: str,
        anchor_name: str,
    ) -> torch.Tensor:
        """
        Compute gradient for a single text span (no paragraph splitting).

        Core computation: backprop cosine similarity w.r.t. hidden state.
        """
        self.model.eval()

        # 1. Get anchor vector (cached, no grad needed)
        anchor_vec = self._get_anchor_vector(anchor_name)

        # 2. Encode article (NEED gradients)
        inputs = self.tokenizer(
            text,
            return_tensors='pt',
            padding=True,
            truncation=True,
            max_length=512,
        )
        inputs = {k: v.to(self.device) for k, v in inputs.items()}

        # Forward pass with hidden states
        output = self.model(**inputs, output_hidden_states=True)

        # Get the layer we want gradients from
        if self.config.gradient_layer == 'last':
            hidden_state = output.last_hidden_state
        elif self.config.gradient_layer == 'second_last':
            hidden_state = output.hidden_states[-2]
        else:  # mean of all layers
            hidden_state = torch.stack(output.hidden_states[1:], dim=0).mean(dim=0)

        # Retain grad so we can inspect it
        hidden_state.retain_grad()

        # Pool
        if self.config.pooling == 'cls':
            article_vec = hidden_state[:, 0, :]
        elif self.config.pooling == 'mean':
            mask = inputs['attention_mask'].unsqueeze(-1).float()
            article_vec = (hidden_state * mask).sum(dim=1) / mask.sum(dim=1)
        else:
            article_vec = hidden_state[:, -1, :]

        # Normalize
        article_vec_norm = F.normalize(article_vec, p=2, dim=1)

        # 3. Compute geometric similarity
        similarity = F.cosine_similarity(article_vec_norm, anchor_vec, dim=1)

        # 4. Backward - gradient of similarity w.r.t. hidden state
        similarity.backward()

        # 5. Extract gradient at pooled position
        if self.config.pooling == 'cls':
            gradient = hidden_state.grad[:, 0, :].detach()
        elif self.config.pooling == 'mean':
            mask = inputs['attention_mask'].unsqueeze(-1).float()
            gradient = (hidden_state.grad * mask).sum(dim=1).detach() / mask.sum(dim=1)
        else:
            gradient = hidden_state.grad[:, -1, :].detach()

        # Optionally normalize gradient
        if self.config.normalize_gradients:
            gradient = F.normalize(gradient, p=2, dim=1)

        # Clean up
        self.model.zero_grad()

        return gradient.squeeze(0)  # [hidden_dim]
    
    def get_all_gradients(
        self, 
        article_text: str,
        anchor_names: List[str] = None,
    ) -> Dict[str, torch.Tensor]:
        """
        Get gradients toward all specified anchors.
        
        Args:
            article_text: The article to analyze
            anchor_names: List of anchor names (default: all)
            
        Returns:
            Dict of {anchor_name: gradient_vector}
        """
        if anchor_names is None:
            anchor_names = list(self.config.anchors.keys())
        
        gradients = {}
        for anchor_name in anchor_names:
            gradients[anchor_name] = self.get_metric_gradient(article_text, anchor_name)
        
        return gradients
    
    def compute_tension(
        self,
        article_text: str,
        anchor_a: str,
        anchor_b: str,
    ) -> Dict[str, float]:
        """
        Compute semantic tension between two framing concepts for an article.
        
        Tension = angle between force vectors toward different anchors.
        
        - High positive correlation: Article pulled equally toward both (ambiguous)
        - Low/negative correlation: Article pulled toward one, away from other (framed)
        
        Returns:
            Dict with correlation, angle, and interpretation
        """
        grad_a = self.get_metric_gradient(article_text, anchor_a)
        grad_b = self.get_metric_gradient(article_text, anchor_b)
        
        # Cosine similarity between gradients
        correlation = F.cosine_similarity(
            grad_a.unsqueeze(0), 
            grad_b.unsqueeze(0), 
            dim=1
        ).item()
        
        # Angle in degrees
        angle = np.arccos(np.clip(correlation, -1, 1)) * 180 / np.pi
        
        # Interpretation
        if correlation > 0.7:
            interpretation = "parallel_pulls"  # Ambiguous, equally pulled
        elif correlation > 0.3:
            interpretation = "mild_tension"
        elif correlation > -0.3:
            interpretation = "orthogonal"  # Independent framings
        elif correlation > -0.7:
            interpretation = "opposing_tension"
        else:
            interpretation = "strong_opposition"  # Locked into one frame
        
        return {
            'anchor_a': anchor_a,
            'anchor_b': anchor_b,
            'correlation': correlation,
            'angle_degrees': angle,
            'interpretation': interpretation,
        }
    
    def compute_tension_matrix(
        self,
        article_text: str,
        anchor_names: List[str] = None,
    ) -> Dict[str, Any]:
        """
        Compute full tension matrix for all anchor pairs.
        
        Returns:
            Dict with matrix, labels, and summary stats
        """
        if anchor_names is None:
            anchor_names = list(self.config.anchors.keys())
        
        n = len(anchor_names)
        gradients = self.get_all_gradients(article_text, anchor_names)
        
        # Stack gradients
        grad_stack = torch.stack([gradients[a] for a in anchor_names], dim=0)  # [n, hidden]
        
        # Compute pairwise cosine similarities
        grad_norm = F.normalize(grad_stack, p=2, dim=1)
        tension_matrix = torch.mm(grad_norm, grad_norm.t())  # [n, n]
        
        # Find most opposing pair
        triu_mask = torch.triu(torch.ones(n, n), diagonal=1).bool()
        upper_vals = tension_matrix[triu_mask]
        min_idx = upper_vals.argmin()
        
        # Convert flat index back to (i, j)
        idx_pairs = [(i, j) for i in range(n) for j in range(i+1, n)]
        most_opposing = idx_pairs[min_idx.item()]
        
        return {
            'matrix': tension_matrix.cpu().numpy(),
            'labels': anchor_names,
            'most_opposing_pair': (anchor_names[most_opposing[0]], anchor_names[most_opposing[1]]),
            'most_opposing_correlation': upper_vals[min_idx].item(),
            'mean_correlation': upper_vals.mean().item(),
            'std_correlation': upper_vals.std().item(),
        }

    def compute_structural_antagonism(
        self,
        article_text: str,
        hypothesis_anchor: str,
        null_anchor: str,
    ) -> Dict[str, Any]:
        """
        Structural Antagonism (Holographic Truth Protocol, Axiom 4).

        Instead of measuring gradient MAGNITUDE (which conflates bias with
        confusion), we measure the PHASE relationship between the gradient
        toward a hypothesis H and the gradient toward its negation ¬H.

            S = -cos(∇L(H), ∇L(¬H))

        Interpretation:
            S ≈ +1  → Anti-parallel gradients → genuine bias (article is
                      pulled TOWARD H and AWAY from ¬H)
            S ≈  0  → Orthogonal gradients → confused / ambiguous
            S ≈ -1  → Parallel gradients → model can't distinguish H from ¬H
                      (degenerate, both attracting equally)

        Args:
            article_text: The article to analyze
            hypothesis_anchor: Name of the hypothesis anchor (e.g., "pro_israel")
            null_anchor: Name of the counter-hypothesis anchor (e.g., "pro_palestine")

        Returns:
            Dict with antagonism score, raw gradients, and interpretation
        """
        grad_h = self.get_metric_gradient(article_text, hypothesis_anchor)
        grad_not_h = self.get_metric_gradient(article_text, null_anchor)

        # S = -cos(∇L(H), ∇L(¬H))
        cos_sim = F.cosine_similarity(
            grad_h.unsqueeze(0),
            grad_not_h.unsqueeze(0),
            dim=1,
        ).item()
        antagonism = -cos_sim

        # Magnitudes (diagnostic only — not used for scoring)
        mag_h = grad_h.norm().item()
        mag_not_h = grad_not_h.norm().item()
        mag_ratio = mag_h / (mag_not_h + 1e-9)

        # Interpretation
        if antagonism > 0.7:
            interpretation = "strong_bias"
        elif antagonism > 0.3:
            interpretation = "moderate_bias"
        elif antagonism > -0.3:
            interpretation = "ambiguous"
        elif antagonism > -0.7:
            interpretation = "weak_signal"
        else:
            interpretation = "degenerate"

        return {
            'hypothesis': hypothesis_anchor,
            'null': null_anchor,
            'antagonism': antagonism,
            'raw_cosine': cos_sim,
            'magnitude_h': mag_h,
            'magnitude_null': mag_not_h,
            'magnitude_ratio': mag_ratio,
            'interpretation': interpretation,
        }

    def compute_antagonism_matrix(
        self,
        article_text: str,
        anchor_pairs: Optional[List[Tuple[str, str]]] = None,
    ) -> Dict[str, Any]:
        """
        Compute structural antagonism for all hypothesis/null pairs.

        If anchor_pairs is None, treats every pair (a, b) with a≠b as
        a candidate hypothesis/null pair.

        Returns:
            Dict with per-pair antagonism scores and summary statistics
        """
        if anchor_pairs is None:
            names = list(self.config.anchors.keys())
            anchor_pairs = [(a, b) for a in names for b in names if a != b]

        results = []
        for h, nh in anchor_pairs:
            results.append(self.compute_structural_antagonism(article_text, h, nh))

        antagonisms = [r['antagonism'] for r in results]
        return {
            'pairs': results,
            'mean_antagonism': float(np.mean(antagonisms)) if antagonisms else 0.0,
            'max_antagonism': float(np.max(antagonisms)) if antagonisms else 0.0,
            'n_biased': sum(1 for a in antagonisms if a > 0.3),
            'n_ambiguous': sum(1 for a in antagonisms if -0.3 <= a <= 0.3),
            'n_degenerate': sum(1 for a in antagonisms if a < -0.7),
        }


# =============================================================================
# BATCH PROCESSING FOR CORPUS ANALYSIS
# =============================================================================

class MetricGradientAnalyzer:
    """
    Analyze entire corpus using metric gradients.
    
    Computes:
    - Per-article tension profiles
    - Cross-article gradient variance
    - Framing clustering based on gradient directions
    """
    
    def __init__(self, extractor: MetricGradientExtractor = None):
        self.extractor = extractor or MetricGradientExtractor()
    
    def analyze_corpus(
        self,
        articles: List[str],
        anchor_pairs: List[Tuple[str, str]] = None,
        article_ids: List[str] = None,
        verbose: bool = True,
    ) -> Dict[str, Any]:
        """
        Analyze corpus for semantic tension patterns.
        
        Args:
            articles: List of article texts
            anchor_pairs: Pairs of anchors to compute tension for
            article_ids: Optional IDs for articles
            
        Returns:
            Comprehensive analysis results
        """
        import datetime
        
        if anchor_pairs is None:
            # Default: test victim vs aggressor
            anchor_pairs = [
                ('victim', 'aggressor'),
                ('emotional', 'neutral'),
                ('humanitarian', 'security'),
            ]
        
        if article_ids is None:
            article_ids = [f"article_{i}" for i in range(len(articles))]
        
        n_articles = len(articles)
        n_pairs = len(anchor_pairs)
        
        if verbose:
            print(f"\n{'='*60}")
            print(f"METRIC GRADIENT CORPUS ANALYSIS")
            print(f"{'='*60}")
            print(f"Articles: {n_articles}")
            print(f"Anchor pairs: {anchor_pairs}")
        
        # Storage
        all_tensions = {f"{a}_{b}": [] for a, b in anchor_pairs}
        all_gradients = {anchor: [] for pair in anchor_pairs for anchor in pair}
        per_article_results = []
        
        for i, (article, aid) in enumerate(zip(articles, article_ids)):
            if verbose and (i + 1) % 10 == 0:
                print(f"  Processing article {i+1}/{n_articles}...")
            
            article_result = {'id': aid, 'tensions': {}}
            
            # Get gradients for all relevant anchors
            relevant_anchors = list(set(a for pair in anchor_pairs for a in pair))
            gradients = self.extractor.get_all_gradients(article, relevant_anchors)
            
            for anchor, grad in gradients.items():
                all_gradients[anchor].append(grad.cpu())
            
            # Compute tensions
            for anchor_a, anchor_b in anchor_pairs:
                grad_a = gradients[anchor_a]
                grad_b = gradients[anchor_b]
                
                corr = F.cosine_similarity(
                    grad_a.unsqueeze(0),
                    grad_b.unsqueeze(0),
                    dim=1
                ).item()
                
                pair_key = f"{anchor_a}_{anchor_b}"
                all_tensions[pair_key].append(corr)
                article_result['tensions'][pair_key] = corr
            
            per_article_results.append(article_result)
        
        # Aggregate statistics
        tension_stats = {}
        for pair_key, tensions in all_tensions.items():
            tensions_arr = np.array(tensions)
            tension_stats[pair_key] = {
                'mean': float(tensions_arr.mean()),
                'std': float(tensions_arr.std()),
                'min': float(tensions_arr.min()),
                'max': float(tensions_arr.max()),
                'median': float(np.median(tensions_arr)),
            }
        
        # Gradient variance per anchor
        gradient_variance = {}
        for anchor, grads in all_gradients.items():
            if grads:
                stacked = torch.stack(grads, dim=0)  # [n_articles, hidden]
                gradient_variance[anchor] = {
                    'variance': float(stacked.var().item()),
                    'mean_norm': float(stacked.norm(dim=1).mean().item()),
                }
        
        result = {
            'n_articles': n_articles,
            'anchor_pairs': anchor_pairs,
            'tension_stats': tension_stats,
            'gradient_variance': gradient_variance,
            'per_article': per_article_results,
            'timestamp': datetime.datetime.now().isoformat(),
        }
        
        if verbose:
            print(f"\n{'='*60}")
            print("ANALYSIS COMPLETE")
            print(f"{'='*60}")
            for pair_key, stats in tension_stats.items():
                print(f"  {pair_key}: mean={stats['mean']:.3f}, std={stats['std']:.3f}")
        
        return result


# =============================================================================
# CONTROL VALIDATION
# =============================================================================

def run_gradient_controls(
    extractor: MetricGradientExtractor,
    real_articles: List[str],
    n_samples: int = 100,
    verbose: bool = True,
) -> Dict[str, Any]:
    """
    Run control experiments for metric gradients.
    
    Tests:
    - Real articles: Should have meaningful gradient variance
    - Constant (same article): Should have ~zero variance
    - Random tokens: Critical test - should NOT have high variance
    
    Expected: real_var >> random_var > constant_var
    """
    import random
    import string
    
    if verbose:
        print(f"\n{'='*60}")
        print("METRIC GRADIENT CONTROL VALIDATION")
        print(f"{'='*60}")
    
    anchor_pair = ('victim', 'aggressor')
    
    # 1. Real articles
    if verbose:
        print(f"\n[1/3] Real articles ({len(real_articles)} samples)...")
    
    real_tensions = []
    for article in real_articles[:n_samples]:
        tension = extractor.compute_tension(article, *anchor_pair)
        real_tensions.append(tension['correlation'])
    
    real_var = np.var(real_tensions)
    real_mean = np.mean(real_tensions)
    
    # 2. Constant (same article repeated)
    if verbose:
        print(f"[2/3] Constant corpus (same article x {n_samples})...")
    
    constant_article = real_articles[0] if real_articles else "This is a test article."
    constant_tensions = []
    for _ in range(n_samples):
        tension = extractor.compute_tension(constant_article, *anchor_pair)
        constant_tensions.append(tension['correlation'])
    
    constant_var = np.var(constant_tensions)
    constant_mean = np.mean(constant_tensions)
    
    # 3. Random tokens
    if verbose:
        print(f"[3/3] Random corpus ({n_samples} samples)...")
    
    def random_text(length=100):
        words = [''.join(random.choices(string.ascii_lowercase, k=random.randint(3, 10))) 
                 for _ in range(length)]
        return ' '.join(words)
    
    random_tensions = []
    for _ in range(n_samples):
        random_article = random_text()
        tension = extractor.compute_tension(random_article, *anchor_pair)
        random_tensions.append(tension['correlation'])
    
    random_var = np.var(random_tensions)
    random_mean = np.mean(random_tensions)
    
    # Analysis
    ordering_satisfied = real_var > random_var > constant_var
    
    if ordering_satisfied:
        verdict = "PASS: Real >> Random >> Constant"
    elif real_var > constant_var:
        verdict = "PARTIAL: Real > Constant, but random ordering unclear"
    else:
        verdict = "FAIL: Expected ordering not satisfied"
    
    result = {
        'real': {'variance': real_var, 'mean': real_mean, 'n': len(real_tensions)},
        'constant': {'variance': constant_var, 'mean': constant_mean, 'n': n_samples},
        'random': {'variance': random_var, 'mean': random_mean, 'n': n_samples},
        'ordering_satisfied': ordering_satisfied,
        'verdict': verdict,
        'anchor_pair': anchor_pair,
    }
    
    if verbose:
        print(f"\n{'='*60}")
        print("CONTROL RESULTS")
        print(f"{'='*60}")
        print(f"  Real variance:     {real_var:.6f} (mean: {real_mean:.3f})")
        print(f"  Constant variance: {constant_var:.6f} (mean: {constant_mean:.3f})")
        print(f"  Random variance:   {random_var:.6f} (mean: {random_mean:.3f})")
        print(f"\n  Verdict: {verdict}")
    
    return result


# =============================================================================
# PLOTLY VISUALIZATION HELPERS
# =============================================================================

def prepare_tension_heatmap_data(analysis: Dict) -> Dict:
    """Prepare data for Plotly heatmap of tension matrix."""
    # This would use the full tension matrix from compute_tension_matrix
    return {
        'note': 'Use extractor.compute_tension_matrix() for single article heatmap',
    }


def prepare_gradient_scatter_data(
    analysis: Dict,
    anchor_pair: Tuple[str, str] = ('victim', 'aggressor'),
) -> Dict:
    """
    Prepare data for Plotly scatter plot of gradient tensions.
    
    X-axis: Tension toward anchor_a
    Y-axis: Tension toward anchor_b
    Each point: One article
    """
    pair_key = f"{anchor_pair[0]}_{anchor_pair[1]}"
    
    tensions = [r['tensions'].get(pair_key, 0) for r in analysis['per_article']]
    ids = [r['id'] for r in analysis['per_article']]
    
    return {
        'x': tensions,
        'y': [1 - abs(t) for t in tensions],  # "Ambiguity" score
        'ids': ids,
        'xlabel': f'Tension ({pair_key})',
        'ylabel': 'Ambiguity (1 - |tension|)',
        'title': f'Semantic Tension Distribution: {pair_key}',
    }


# =============================================================================
# INTEGRATION WITH EMBEDDING PIPELINE
# =============================================================================

def compare_gradient_vs_embedding(
    articles: List[str],
    embedding_results: Dict,  # From DirichletFusion
    gradient_results: Dict,   # From MetricGradientAnalyzer
) -> Dict:
    """
    Compare metric gradient approach with embedding approach.
    
    Answers: Do they find the same structure?
    """
    # Extract per-article measures from both
    n = len(articles)
    
    # From embeddings: use observer variance per article
    emb_variance = embedding_results.get('variance_decomposition', {}).get('v_fraction', 0)
    
    # From gradients: use tension std per article
    grad_tensions = [r['tensions'] for r in gradient_results['per_article']]
    
    # Correlation between approaches would go here
    # This is a placeholder for the actual comparison
    
    return {
        'embedding_v_variance': emb_variance,
        'gradient_tension_stats': gradient_results['tension_stats'],
        'note': 'Full comparison requires aligned per-article metrics',
    }


# =============================================================================
# MODULE EXPORTS
# =============================================================================

__all__ = [
    # Track 3: Geometric Stability (Added for Modern System)
    'compute_metric_gradients',
    'compute_alpha_gradients',
    'AlphaStabilityAnalyzer',
    'AlphaStabilityConfig',
    
    # Track 1.5: Semantic Axis Tension (Preserved Original)
    'MetricGradientConfig',
    'MetricGradientExtractor',
    'MetricGradientAnalyzer',
    'run_gradient_controls',
    'prepare_tension_heatmap_data',
    'prepare_gradient_scatter_data',
    'compare_gradient_vs_embedding',
]


# =============================================================================
# TEST
# =============================================================================

if __name__ == "__main__":
    print("=" * 60)
    print("Testing Metric Gradient Extraction (Dual-Track)")
    print("=" * 60)
    
    # 1. Test Track 1.5 (Anchors)
    config = MetricGradientConfig(
        device='cpu',  # Use CPU for testing
        anchors={
            'victim': "This describes victims and suffering.",
            'aggressor': "This describes aggression and violence.",
        }
    )
    extractor = MetricGradientExtractor(config)
    article = "Israel struck a hospital in Gaza, killing 50 civilians including children."
    
    print(f"\n[Track 1.5] Test article: {article[:60]}...")
    tension = extractor.compute_tension(article, 'victim', 'aggressor')
    print(f"Correlation: {tension['correlation']:.4f}")
    print(f"Interpretation: {tension['interpretation']}")
    
    # 2. Test Track 3 (Stability - Mock)
    print("\n[Track 3] Testing Stability Calculation (Mock Data)...")
    mock_sweep = {
        'wavelength': {
            '0.1_to_0.5': {'energy': 0.05},
            '0.5_to_1.0': {'energy': 0.02},
            '1.0_to_5.0': {'energy': 0.01}
        }
    }
    alphas = [0.1, 0.5, 1.0, 5.0]
    stability = compute_metric_gradients(mock_sweep, alphas)
    print(f"Stability Score: {stability['stability_score']:.4f}")
    
    print("\n" + "=" * 60)
    print("TEST COMPLETE")
    print("=" * 60)
\end{CodeBlock}
\end{document}
